% Part: proof-theory
% Chapter: sequent-calculus
% Section: interpretation-rules

\documentclass[../../../include/open-logic-section]{subfiles}

\begin{document}

\olfileid{pt}{seq}{irl}

\olsection{નિયમોનું અર્થઘટન}

સિક્વન્ટનું અર્થઘટન યાદ કરો:
$\Gamma \Sequent \Delta$ કોઈ !!a{structure}~$\Struct{M}$માં સત્ય છે,
જો ઓછામાં ઓછું એક $!A \in \Gamma$ અસત્ય હોય
અથવા ઓછામાં ઓછું એક~$!A \in \Delta$ સત્ય હોય.
\Log{G3c}ના તાર્કિક નિયમોને તેમના મુખ્ય !!{formula}sની
સત્યતાની શરતો રજૂ કરતા નિયમો તરીકે વાંચી શકીએ.
\RightR{\lif} વિચારીએ. તેનું મુખ્ય !!{formula} $!A \lif !B$ છે.
$!A$ અસત્ય હોય અથવા~$!B$ સત્ય હોય, તો તે સત્ય છે.
એટલે $\ \Sequent !A \lif !B$ સત્ય છે
જો અને ફક્ત જો $!A \Sequent !B$ સત્ય હોય.
આ સિક્વન્ટોના પૂર્વાંગ અને ઉત્તરાંગમાં સંદર્ભનાં
!!{formula}s $\Gamma$, $\Delta$ ઉમેરીએ,
તો બરાબર \RightR{\lif} નિયમનો નિષ્કર્ષ અને આધાર મળે છે.
\LeftR{\lif} માટે પણ સ્થિતિ સમાન છે.
$!A$~સત્ય અને $!B$~અસત્ય હોય, તો $!A \lif !B$ અસત્ય છે.
એટલે $!A \lif !B \Sequent \ $ સત્ય છે
જો અને ફક્ત જો બંને $\ \Sequent !A$ અને $!B \Sequent \ $ સત્ય હોય.
\LeftR{\lif}નો નિષ્કર્ષ તેના આધારોના સંયોજનને સમકક્ષ છે.
આ સ્વરૂપ~\Log{G3c}ના નિયમોની યથાર્થતા સુનિશ્ચિત કરે છે
(બધા આધાર સત્ય હોય, તો નિષ્કર્ષ સત્ય છે).
તે પૂર્ણતા પણ સુનિશ્ચિત કરે છે.

કોઈ પણ સત્ય-વિધેયાત્મક સંયોજકની સત્યતાની શરતો
આ રીતે સિક્વન્ટોના સમૂહથી વર્ણવી શકાય.
આ એ હકીકતમાંથી અનુસરે છે કે શાસ્ત્રીય તર્કમાં
દરેક સત્ય-વિધેયને સંયોજક સામાન્ય સ્વરૂપના !!a{formula}થી વર્ણવી શકાય:
આણ્વિક અને નિષેધિત આણ્વિક !!{formula}sના વિકલ્પોનું સંયોજન.
દાખલા તરીકે $!A \oplus !B$ને સત્ય ગણીએ
જો અને ફક્ત જો $!A$ અને $!B$માંથી બરાબર એક સત્ય હોય,
એટલે ``અપવર્જક વિકલ્પ.''
ત્યારે $!A \oplus !B$ સત્ય છે જો અને ફક્ત જો
$(!A \lor !B) \land (\lnot !A \lor \lnot !B)$ સત્ય હોય,
અને $(\lnot !A \lor !B) \land (!A \lor \lnot !B)$ સત્ય હોય
તો તે અસત્ય છે. એટલે $\oplus$ માટે
સિક્વન્ટ કલનના નિયમો આ રીતે રજૂ કરી શકીએ:
\begin{align*}
&
\Axiom$\Gamma \fCenter \Delta, !A, !B$
\Axiom$!A, !B, \Gamma \fCenter \Delta$
\RightLabel{\RightR{\oplus}}
\BinaryInf$\Gamma \fCenter \Delta, !A \oplus !B$
\DisplayProof
\\
& \Axiom$!A, \Gamma \fCenter \Delta, !B$
\Axiom$!B, \Gamma \fCenter \Delta, !A$
\RightLabel{\LeftR{\oplus}}
\BinaryInf$!A \oplus !B, \Gamma \fCenter \Delta$
\DisplayProof
\end{align*}
આ નિયમો પણ યથાર્થ અને પૂર્ણ હોય.
તેમનાં નિષ્કર્ષ સત્ય છે જો અને ફક્ત જો બધા આધાર સત્ય હોય.
\sourcecorrection{OLMD-304}{મૂળના બીજા અપવર્જક-વિકલ્પ નિયમમાં મુખ્ય સૂત્ર ડાબે છે પરંતુ નિયમ-લેબલ જમણો કહે છે. તેને ડાબો નિયમ કર્યો છે; આધાર અને નિષ્કર્ષ અચળ છે.}

\begin{prob}
$!A \downarrow !B$ સત્ય છે જો અને ફક્ત જો
$!A$ કે~$!B$માંથી કોઈ સત્ય ન હોય (``NOR'') એમ ધારો.
તેના \Log{G3c} નિયમો કયા હોય?
(એવા બે સિક્વન્ટો શોધો જે એકસાથે સત્ય હોય
જો અને ફક્ત જો $!A \downarrow !B$ સત્ય હોય.
તેથી \RightR{\downarrow} નિયમ મળે.
પછી એવો એક સિક્વન્ટ શોધો જે સત્ય હોય
જો અને ફક્ત જો $!A \downarrow !B$ અસત્ય હોય.
તેથી \LeftR{\downarrow} નિયમ મળે.)
\end{prob}

\Log{G3c}માં દરેક સંયોજક અને પરિમાણકના બરાબર બે નિયમ છે:
એક ડાબો અને એક જમણો.
પરંતુ \Log{G1c}માં બે \LeftR{\land} અને બે \RightR{\lor} નિયમ છે.
કારણ એ છે કે \Log{G1c} જેનું અનુસરણ કરે છે તે
ગેન્ટ્ઝેનના મૂળ તંત્ર~\Log{LK}નું,
મહત્ત્વના અશાસ્ત્રીય તર્ક એટલે અંતઃપ્રજ્ઞાવાદી તર્ક માટે,
\Log{LJ} નામનું સ્વરૂપ હતું.
અંતઃપ્રજ્ઞાવાદી તર્ક માટે યથાર્થ અને પૂર્ણ તંત્ર મેળવવા
\Log{LJ}ના દરેક સિક્વન્ટના ઉત્તરાંગમાં
વધુમાં વધુ એક !!{formula} રાખવાની મર્યાદા છે.
એટલે~\Log{G3c}નો \RightR{\lor} નિયમ વાપરી ન શકાય,
કારણ કે તેના આધારના ઉત્તરાંગમાં બંને $!A$ અને~$!B$ છે.
તેથી ગેન્ટ્ઝેને \Log{LJ} માટે \RightR{\lor}ના
બે નિયમો વાપર્યા અને \Log{LK} માટે પણ એ જ વાપર્યા.
અંતઃપ્રજ્ઞાવાદી સ્વરૂપ~\Log{G1i} એ~\Log{G1c}ને અનુરૂપ છે,
અને તેમાં બધા ઉત્તરાંગો વધુમાં વધુ એક !!{formula} ધરાવે છે.
\sourcecorrection{OLMD-305}{મૂળ G1iને ફક્ત ઉત્તરાંગની મર્યાદા ધરાવતું G1c કહે છે. પરંતુ સ્થિર G1i કોષ્ટકના અનુસરણ-ડાબા નિયમના પહેલા આધારમાં મુખ્ય પૂર્વાંગ સિવાયનો ઉત્તરાંગ નથી; બીજો આધાર પોતાની મર્યાદિત ઉત્તરાંગ-સામગ્રી રાખે છે. એટલે ફક્ત દરેક સિક્વન્ટને મર્યાદિત કરવાથી બધા અંતઃપ્રજ્ઞાવાદી નિયમો આપમેળે મળતા નથી; સંબંધિત નિયમનું સ્વરૂપ પણ બદલાય છે.}
(\Log{G3c}નું પણ અંતઃપ્રજ્ઞાવાદી સ્વરૂપ છે,
પરંતુ તેના કેટલાક નિયમો~\Log{G3c}ના અનુરૂપ નિયમોથી જુદા છે.
દાખલા તરીકે તેમાં પણ \RightR{\lor}ના બે નિયમો છે.)

\Log{G1c}ના બંધારણીય નિયમો યથાર્થ છે તે સરળતાથી દેખાય છે.
દાખલા તરીકે $\Gamma \Sequent \Delta$માં
જમણે સત્ય !!{formula} અથવા ડાબે અસત્ય !!{formula} હોય,
તો $\Gamma \Sequent \Delta, !A$માં પણ એ જ સૂત્ર છે;
$!A$ સત્ય છે કે અસત્ય તે અપ્રસ્તુત છે.
એટલે~\RightR{\Weakening}નો આધાર સત્ય હોય, તો નિષ્કર્ષ પણ સત્ય છે.
અહીં ઊલટું સાચું નથી તે નોંધો.
$!A$ સત્ય હોવાથી નિષ્કર્ષ સત્ય હોઈ શકે;
ત્યારે આધાર સત્ય હોવાની ખાતરી નથી.
\sourcecorrection{OLMD-306}{મૂળ આ વાક્યમાં જમણે અસત્ય અને ડાબે સત્ય સૂત્ર કહે છે. તે સિક્વન્ટની સત્યતાની શરતથી ઊલટું છે. જમણે સત્ય અથવા ડાબે અસત્ય કર્યું છે, જેમ વિભાગની શરૂઆતની વ્યાખ્યામાં છે.}

બંધારણીય નિયમો શા માટે જોઈએ?
દાખલા તરીકે $\ \Sequent !A \lor \lnot !A$ને
!!{prove} કરવા સંકોચન (\RightR{\Contraction}, \LeftR{\Contraction}) જોઈએ.
$!A \Sequent !A$થી શરૂ કરીએ,
તો~\RightR{\lnot}થી પહેલા $\Sequent !A, \lnot !A$ મળે.
પછી \RightR{\lor}નાં બંને સ્વરૂપ એક પછી એક વાપરી શકીએ:
એક વાર $\lnot !A$માંથી $!A \lor \lnot !A$ મેળવવા
અને એક વાર~$!A$માંથી $!A \lor \lnot !A$ મેળવવા.
$\ \Sequent !A \lor \lnot !A, !A \lor \lnot !A$ મળે છે.
એટલે~\Log{G1c}માં $\Sequent !A \lor \lnot !A$ મેળવવા
!!a{proof}માં એ જ !!{formula}ની બે આવૃત્તિનું
``સંકોચન'' કરી શકવું જોઈએ:
\[
\Axiom$!A \fCenter !A$
\RightLabel{\RightR{\lnot}}
\UnaryInf$\fCenter !A, \lnot !A$
\RightLabel{\RightR{\lor}}
\UnaryInf$\fCenter !A, !A \lor \lnot !A$
\RightLabel{\RightR{\lor}}
\UnaryInf$\fCenter !A \lor \lnot !A, !A \lor \lnot !A$
\RightLabel{\RightR{\Contraction}}
\UnaryInf$\fCenter !A \lor \lnot !A$
\DisplayProof
\]
ખરેખર કેટલાક માન્ય સિક્વન્ટો~\Log{G1c}માં
શિથિલીકરણ કે સંકોચન વિના !!{prove} કરી ન શકાય.
દાખલા તરીકે~\Log{G1c}માં $!A \Sequent !B \lif !A$
સિક્વન્ટની !!a{proof}નો અંત \RightR{\lif}થી જ થાય
(કારણ કે ફક્ત $!B \lif !A$ મુખ્ય સૂત્ર હોઈ શકે).
તે માટે આધાર તરીકે $!B, !A \Sequent !A$ જોઈએ.
તેને સ્વયંસિદ્ધ $!A \Sequent !A$માંથી મેળવવા
\LeftR{\Weakening} જોઈએ:
\[
\Axiom$!A \fCenter !A$
\RightLabel{\LeftR{\Weakening}}
\UnaryInf$!B, !A \fCenter !A$
\RightLabel{\RightR{\lif}}
\UnaryInf$!A \fCenter !B \lif !A$
\DisplayProof
\]

\Log{G3c}માં સંકોચન અને શિથિલીકરણના નિયમો જરાય જરૂરી નથી.
શિથિલીકરણ સ્વયંસિદ્ધોમાં જ સમાયેલું છે.
બે આધાર ધરાવતા નિયમોમાં સંદર્ભ $\Gamma$, $\Delta$ની
બે નકલને એકમાં ભેળવીને મોટેભાગે સંકોચન ટાળી શકાય.
\Log{G1c} અને \Log{G3c} બંનેમાં આવા
``સંદર્ભ વહેંચતા'' નિયમો છે.
એક આધારવાળા નિયમોમાં ફક્ત \RightR{\lor}, \LeftR{\land},
\RightR{\lexists} અને~\LeftR{\lforall} માટે સંકોચન જોઈએ.
\Log{G3c}માં આ નિયમોનાં સ્વરૂપ સંકોચનને સંપૂર્ણપણે બિનજરૂરી બનાવે છે.
એક \Log{G2c} તંત્ર પણ છે (\olref{tab:G2c} જુઓ),
જેમાં બે આધારવાળા નિયમોના સંદર્ભો સ્વતંત્ર હોય છે.
\Log{G2c}માં સંકોચન~\Log{G1c} કરતાં પણ વધુ મહત્ત્વનું છે.

\subfile{rules-G2c}

\end{document}
